\begin{figure}[H]
    \centering
    \scriptsize
    \setlength{\tabcolsep}{.4em}
    \begin{tabular}{l|l}
        \begin{lstlisting}
\ Complex Fast Fourier Transform
\ Usage: Vector.name{ N FORWARD ( INVERSE ) }FFT

TASK FFT
FIND  C+        0= ?( FLOAD COMPLEX )
FIND  1ARRAY    0= ?( FLOAD MATRIX.HSF )
FIND  FILL      0= ?( FLOAD FILEIO.FTH )
FIND  TRIG      0= ?( FLOAD TRIG )
\ if not there
DECIMAL
\ = = = = = = = = = = = = = = = = = = = = = = = = = =
\ auxiliary words
  CODE  SHR  BX 1 SHR.  END-CODE
  : LG2 ( n -- lg2[n]) 0 SWAP ( -- 0 n)  SHR
    BEGIN   ?DUP 0>
    WHILE   SHR   SWAP 1+  SWAP  REPEAT  ;

0 VAR DIRECTION?
: FORWARD   0 IS DIRECTION? ;
: INVERSE  -1 IS DIRECTION? ;

    0 VAR N.BITS            \ some VARs
    0 VAR N
    0 VAR MMAX
    0 VAR f{

: C/N    N S->F  C/F  ;
: NORMALIZE    DIRECTION?
    IF   C/N  CPSWAP  C/N   CPSWAP   THEN ;
\ end auxiliary words ;
\ = = = = = = = = = = = = = = = = = = = = = = = = = =
\ key bit-reveral routine!
0 VAR I.R
: B.R  ( n -- n')             \ reverses order of bits
     0  SWAP   ( -- 0 n )     \ set up stack
     N.BITS 0 DO  DUP  1 AND  \ pick out 1's bit
            ROT  2* +     \ double sum and add 1's bit
            SWAP  2/      \ n -> n/2
     LOOP  DROP ;

: BIT.REVERSE  0 IS I.R
     N 0  DO   I B.R   IS  I.R
       I.R  I  < NOT   ( I.R >=I ?)
       IF  f{ I.R } G@L  f{ I }  G@L  NORMALIZE
           f{ I.R } G!L  f{ I }  G!L  THEN
     LOOP ;
\ end bit-reversal (N times)
        \end{lstlisting}
      &
        \begin{lstlisting}
\ = = = = = = = = = = = = = = = = = = = = = = = = = =
                                    \ main algorithm
CODE C-+   2 FLD.   1 FXCH.  3 FSUBR".
    1 FADDP.  1 FXCH. 3 FLD.
    1 FXCH.  4 FSUBR".  1 FADDP.  1 FXCH.  END-CODE
    ( 87: w z -- w-z  w+z)

: THETA   F=PI  MMAX  S->F  F/  DIRECTION? FSIGN ;

CREATE  WP  16  ALLOT OKLW
: INIT.TRIG   FINIT  THETA  EXP(I*PHI)  WP  DCP!  C=1 ;

: NEW.W   ( 87: w -- w')  WP DCP@ C* ;

0 VAR ISTEP

: DO.INNER.LOOP
    DO  MMAX I +  IS I.R
        CPDUP  f{ I.R } G@L  C*   f{ I } G@L
        CPSWAP  C-+ f{ I } G!L  f{ I.R } G!L
    ISTEP +LOOP ;

: }FFT ( adr n -- ) IS N IS f{
    1 IS MMAX
    N LG2 IS N.BITS
    FINIT
    BIT.REVERSE
    BEGIN
      N MMAX >
    WHILE
      INIT.TRIG  MMAX  2* IS ISTEP
      MMAX 0 DO
        N I  DO.INNER.LOOP   NEW.W
      LOOP
      ISTEP IS MMAX
    REPEAT  CPDROP  ;

: POWER  0 DO  f{ I } G@L  CABS  CR  I .  F.  LOOP ;
\ power spectrum of FFT            \ end of fft code
\ = = = = = = = = = = = = = = = = = = = = = = = = = =
                                   \ an example
64 LONG COMPLEX 1ARRAY A{
: INIT.A  A{ $" IBM.EX" OPEN-INPUT FILL CLOSE-INPUT  ;
INIT.A
A{ 64 DIRECT }FFT
                                   \ end of example
\ = = = = = = = = = = = = = = = = = = = = = = = = = =
        \end{lstlisting}
    \end{tabular}
\end{figure}